\begin{proposition}[Full cyclotomic action on the Frattini module]
\label{prop:full-cyclotomic-frattini-action}
Retain \(P=M/R_M\), \(V=P/\Phi(P)\), and the coordinate
\(V=\mathbf F_2[[T]]\), where \(T=h_*-1\) and
\(\lambda_2(h_*)=-3\).
Put
\[
 U=G/M=\widehat{\mathbf Z}^{\,\times},\qquad
 \Gamma=1+4\mathbf Z_2,\qquad
 U_0=\{\pm1\}\times\prod_{q\ {\rm odd}}\mathbf Z_q^\times.
\]
Thus \(U=\Gamma\times U_0\).
For \(u\in U\), let \(\langle u_2\rangle\) be the unique member
of \(\{u_2,-u_2\}\) in \(\Gamma\).
The complete action on \(V\) is
\begin{equation}
 u\cdot F(T)
   =(1+T)^{\,\log\langle u_2\rangle/\log(-3)}F(T).        \label{eq:cyclotomic-frattini-1}
\end{equation}
In particular \(U_0\) acts trivially.
\end{proposition}

\begin{proof}
The group \(M\) acts on \(P\) by inner automorphisms and hence
trivially on \(V\). The action factors through the abelian group
\(U\), so every action operator commutes with \(h_*\), and hence
with \(T\). By continuity it is an endomorphism of the cyclic
completed module \(\mathbf F_2[[T]]\).
An automorphism is therefore multiplication by a unit.
This gives a continuous homomorphism
\[
 U\longrightarrow\mathbf F_2[[T]]^\times;
\]
continuity follows also by applying the operator to the cyclic
vector \(1=[a_3]\).

The target is a torsion-free pro-two group.
Indeed it is \(1+T\mathbf F_2[[T]]\), and its quotients modulo
powers of \(T\) are finite two-groups.
If \(a^{2^n}=1\), then \((a-1)^{2^n}=0\); the power-series
domain forces \(a=1\).
Consequently the finite dyadic sign factor acts trivially.
For odd \(q\), the open pro-\(q\) subgroup of
\(\mathbf Z_q^\times\) has trivial image in a pro-two group,
and the remaining finite image is trivial by torsion-freeness.
Finite-support products are dense in the odd-local product,
so the whole of \(U_0\) acts trivially.

It remains to determine the action of \(\Gamma\).
Its topological generator \(-3\) acts as \(1+T\) by the exact
Frattini recursion. This gives \eqref{eq:cyclotomic-frattini-1}.
The exponent lies in \(\mathbf Z_2\), since \(-3\)
topologically generates \(1+4\mathbf Z_2\).
\end{proof}

\begin{proposition}[The residual extension obstruction]
\label{prop:explicit-residual-extension-obstruction}
Let \(W=H^1(P,\mathbf F_2)\), with its discrete continuous
\(U\)-action. There is a natural isomorphism
\begin{equation}
 H^1_{\rm cts}(U,W)
   \cong\operatorname{Hom}_{\rm cts}(U_0,\mathbf F_2).    \label{eq:cyclotomic-frattini-2}
\end{equation}
In particular the latter is a direct sum of one copy of
\(\mathbf F_2\) for the dyadic sign factor and one for each
odd rational prime.
There is an exact segment
\begin{equation}
 0\longrightarrow\mathbf F_2(\mu\bmod2)
 \longrightarrow H^1(M,\mathbf F_2)^G
 \longrightarrow H^1(R_M,\mathbf F_2)^G
 \longrightarrow
       \operatorname{Hom}_{\rm cts}(U_0,\mathbf F_2),     \label{eq:cyclotomic-frattini-3}
\end{equation}
where \(\mu=\rho_{2,2}^I/2\).
No surjectivity of the last map is asserted.
\end{proposition}

\begin{proof}
Let \(\ell_r\) be dual to \(T^r\), so
\(W=\bigoplus_{r\ge0}\mathbf F_2\ell_r\).
For the pullback action of the chosen generator,
\[
 \ell_r\circ(h_*-1)=\ell_{r-1}\quad(r\ge1),\qquad
 \ell_0\circ(h_*-1)=0.
\]
The inverse-pullback convention has the same image and kernel.
Thus generator minus one acts surjectively on \(W\), and
\(W^\Gamma=\mathbf F_2\ell_0\).

A continuous cocycle on the procyclic group \(\Gamma\) is
determined by its value on a topological generator, and every
value in \(W\) is permitted.
Indeed a vector has finite orbit of two-power size; summing
around that orbit twice gives zero, so its cocycle factors
through a sufficiently large finite cyclic quotient.
Coboundaries are the image of generator minus one.
It follows that \(H^1_{\rm cts}(\Gamma,W)=0\).

Given a \(U\)-cocycle \(z\), subtract a coboundary so that its
restriction to \(\Gamma\) is zero.
Since \(U=\Gamma\times U_0\), comparison of the cocycle identities
for \(\gamma u=u\gamma\) gives \(z(u)\in W^\Gamma\) for
\(u\in U_0\).
The action of \(U_0\) on \(W\) is trivial by the preceding
proposition. Hence this restriction is a continuous
homomorphism \(U_0\to\mathbf F_2\ell_0\).
Conversely every such homomorphism, extended by zero on
\(\Gamma\), is a \(U\)-cocycle.
Changing the normalization by a \(\Gamma\)-invariant
coboundary has no effect on \(U_0\). This proves \eqref{eq:cyclotomic-frattini-2}.

For odd \(q\), the pro-\(q\) unit subgroup contributes no
quadratic character, and the cyclic residue-unit group
contributes one. Continuity permits only finitely many
nonzero components, proving the direct-sum description.

Finally use the \(U\)-equivariant exact sequence
\[
 0\longrightarrow W
 \longrightarrow H^1(M,\mathbf F_2)
 \longrightarrow H^1(R_M,\mathbf F_2)^M
 \longrightarrow0
\]
proved by freeness of \(P\).
Taking invariants and applying \eqref{eq:cyclotomic-frattini-2}, with
\(W^U=\mathbf F_2\ell_0=\mathbf F_2(\mu\bmod2)\),
gives \eqref{eq:cyclotomic-frattini-3}.
\end{proof}

\begin{remark}
Concretely, an invariant character of \(R_M\) first extends to
\(M\). Its conjugation defect is a \(U\)-cocycle in \(W\).
An adjustment by a character of \(P\) removes the
\(\Gamma\)-defect. The remaining defect is a finite-support
quadratic character of \(U_0\), multiplied by the mass
functional. This describes the obstruction to an invariant
extension; it does not make that obstruction or the relative
abelian quotient of \(R_D\) vanish.
\end{remark}
